\section{Tessuto osseo}

\rubrica{Funzioni}
\begin{itemize}
  \item sostegno e supporto strutturale all'intero organismo;
  \item possibilità di traslazione di parti dell'organismo (sistema di leve);
  \item metabolismo di calcio ed altri ioni;
  \item protezione di organi ed apparati;
  \item sostegno e protezione a tessuto emopoietico.
\end{itemize}

% ---------------------------------------------------------------------------
\subsection{Matrice extracellulare}

\rubrica{Componente organica (33\%)}
\begin{itemize}
  \item \textbf{collagene};
  \item \textbf{componente amorfa}: GAG (condroitinsolfato, cheratansolfato e acido ialuronico); glicoproteine (ostenectina e osteoclacina).
\end{itemize}

\rubrica{Componente inorganica (67\%)}
\begin{itemize}
  \item fosfato di calcio (cristalli di idrossiapatite), 39\%;
  \item carbonato (9,8\%) e fosfato (17\%) di calcio;
  \item fluoruro di calcio;
  \item fosfato di magnesio;
  \item altri elementi (K, Na, Mn).
\end{itemize}

% ---------------------------------------------------------------------------
\subsection{Componente cellulare}

Cellule osteoprogenitrici, osteoblasti e osteociti rappresentano tre stadi evolutivi della stessa cellula di origine mesenchimale. Gli osteoclasti, cellule giganti plurinucleate, derivano invece da progenitori della linea monocito-macrofagica.

\rubrica{Cellule progenitrici}
Danno origine agli osteoblasti; cellula staminale la cui divisione produce osteoblasti.

\rubrica{Osteoblasti}
Cellula ossea immatura che produce i componenti organici della matrice.
\begin{itemize}
  \item producono la parte proteica della matrice;
  \item stimolano la osteogenesi;
  \item danno origine agli osteociti;
  \item aspetto globoso, molto voluminosi;
  \item allineati sulla matrice in corso di deposizione;
  \item man mano che vengono racchiusi nelle lacune, si trasformano in osteociti.
\end{itemize}

\rubrica{Osteociti}
Cellula ossea matura che mantiene la matrice dell'osso.
\begin{itemize}
  \item situati in spazi ricavati tra le lamelle (\textbf{lacune ossee});
  \item mantenimento e controllo di proteine e minerale;
  \item partecipazione a riparazione di tessuto osseo danneggiato;
  \item forma ellissoide, appiattita e allungata, con asse maggiore parallelo alla direzione delle fibre;
  \item prolungamenti più o meno lunghi e ramificati che decorrono all'interno dei \textbf{canalicoli};
  \item osteociti contigui comunicano attraverso \textit{gap junctions} presenti sulle parti terminali dei prolungamenti.
\end{itemize}

\rubrica{Osteoclasti}
Cellula multinucleata che secerne acidi ed enzimi che erodono la matrice ossea.
\begin{itemize}
  \item degradano la matrice ossea;
  \item polinucleati e di notevoli dimensioni;
  \item forma indefinita e notevole motilità;
  \item quando attivati, aderiscono alla matrice ossea da riassorbire:
  \begin{enumerate}
    \item \textbf{demineralizzazione}: eliminazione dei sali minerali di idrossiapatite;
    \item eliminazione della sostanza organica proteica della matrice ossea ad opera dei lisosomi.
  \end{enumerate}
\end{itemize}

% ---------------------------------------------------------------------------
\subsection{Periostio ed endostio}

\rubrica{Periostio}
\begin{itemize}
  \item connettivo fibrillare denso a fasci intrecciati;
  \item strato esterno fibroso;
  \item strato interno cellulare e vascolarizzato;
  \item alcune cellule hanno capacità osteoformativa $\rightarrow$ crescita ossea per apposizione;
  \item dallo strato profondo si dipartono robusti fasci di fibre $\rightarrow$ \textbf{fibre perforanti di Sharpey}, che garantiscono l'ancoraggio del connettivo all'osso.
\end{itemize}

\rubrica{Endostio}
\begin{itemize}
  \item riveste il canale midollare;
  \item molto vascolarizzato;
  \item presenti cellule osteogenetiche ed ematopoietiche.
\end{itemize}

% ---------------------------------------------------------------------------
\subsection{Classificazione in base all'organizzazione della matrice ossea}

\begin{center}
\begin{forest} classalbero
[Tessuto osseo
  [Non lamellare
    [A fasci intrecciati]
    [A fasci paralleli]
  ]
  [Lamellare
    [Compatto]
    [Spugnoso]
  ]
]
\end{forest}
\end{center}

\rubrica{Tessuto osseo non lamellare}
Matrice disposta a formare una massa compatta; costituisce l'osso primario dei mammiferi.

\rubrica{Tessuto osseo lamellare}
Matrice suddivisa in strutture lamellari.

% ---------------------------------------------------------------------------
\subsection{Tessuto osseo non lamellare}

\rubrica{A fasci intrecciati}
Detto anche \textbf{osso fibroso}.
\begin{itemize}
  \item presente nelle suture e nelle inserzioni di tendini e legamenti;
  \item lacune ossee distribuite irregolarmente;
  \item fibre collagene disposte irregolarmente secondo fasci intrecciati;
  \item presenti cavità contenenti vasi.
\end{itemize}

\rubrica{A fasci paralleli}
Detto anche \textbf{pseudolamellare}.
\begin{itemize}
  \item presente solo temporaneamente nella prima formazione delle ossa lunghe;
  \item fibre collagene disposte in grossi fasci paralleli;
  \item lacune ossee piccole, orientate con l'asse maggiore parallelo alle fibre collagene.
\end{itemize}

% ---------------------------------------------------------------------------
\subsection{Tessuto osseo lamellare}

Matrice disposta in lamelle, al cui interno sono scavate le lacune ossee contenenti gli osteociti; fibre collagene parallele fra di loro.

\rubrica{Tessuto osseo lamellare compatto}
Lamelle strettamente sovrapposte una all'altra in anelli concentrici (\textbf{osteoni}) o in lamine parallele.

\rubrica{Tessuto osseo lamellare spugnoso}
Lamelle disposte in modo da formare \textbf{trabecole} che presentano numerose cavità.

\rubrica{Localizzazione}
Nell'osso lungo si distinguono epifisi, metafisi e diafisi. L'osso spugnoso è presente dove le forze vengono applicate da varie direzioni (es.\ epifisi); l'osso compatto è molto resistente alla compressione in senso longitudinale, ma una pressione laterale provoca facilmente fratture (es.\ diafisi).

% ---------------------------------------------------------------------------
\subsection{Tessuto osseo lamellare compatto}

\rubrica{Osteone (sistema haversiano)}
Lamelle ossee concentriche disposte attorno ad un canale centrale (\textbf{canale di Havers}), dove sono situati vasi nutritizi. Le lamelle presentano lacune contenenti gli osteociti.

\rubrica{Lamelle concentriche}
\begin{itemize}
  \item si dispongono attorno al canale di Havers (da 8 a 15 lamelle);
  \item quella più interna è la più recente;
  \item la più esterna è delimitata da una linea frastagliata più mineralizzata, detta \textbf{linea cementante};
  \item lamelle strettamente aderenti tra loro, unite da una matrice molto mineralizzata;
  \item spazi tra le lamelle assenti o molto ridotti;
  \item le fibre collagene decorrono parallele con andamento elicoidale, ruotate di 90° rispetto alle fibre della lamella adiacente.
\end{itemize}

\rubrica{Lamelle e canalicoli}
\begin{itemize}
  \item le lacune ossee contenenti gli osteociti sono scavate nello spessore delle lamelle, con asse maggiore parallelo alle fibre collagene;
  \item da ogni lacuna si dipartono dei canalicoli che mettono in comunicazione le varie lamelle, garantendo gli scambi metabolici tra gli osteociti.
\end{itemize}

\rubrica{Microanatomia della diafisi}
Nella diafisi di un osso lungo gli osteoni sono disposti parallelamente tra loro e parallelamente all'asse principale dell'osso.
\begin{itemize}
  \item \textbf{lamelle circonferenziali esterne ed interne}: rivestono, rispettivamente, la superficie esterna ed interna della diafisi;
  \item \textbf{lamelle interstiziali}: residui di osteoni preesistenti, interposte tra gli osteoni attuali;
  \item il \textbf{canale di Havers} decorre parallelo all'asse dell'osso; il \textbf{canale di Volkmann} lo attraversa trasversalmente, mettendo in comunicazione i canali di Havers tra loro e con periostio ed endostio;
  \item è presente una piccola quantità di tessuto osseo spugnoso, con trabecole e cavità contenenti midollo osseo, in comunicazione con la cavità centrale della diafisi.
\end{itemize}

% ---------------------------------------------------------------------------
\subsection{Tessuto osseo lamellare spugnoso}

\begin{itemize}
  \item presente nella parte interna delle ossa brevi e nelle epifisi delle ossa lunghe;
  \item presenza di un trabecolato che dà all'osso un aspetto spugnoso;
  \item presenza, nelle cavità, di midollo osseo, vasi e nervi.
\end{itemize}
Le trabecole ossee, più o meno spesse, delimitano cavità midollari comunicanti contenenti midollo osseo, vasi e nervi; al loro interno sono visibili adipociti e lacune ossee che accolgono gli osteociti.

% ---------------------------------------------------------------------------
\subsection{Ossificazione}

\begin{center}
\begin{forest} classalbero
[Ossificazione
  [Diretta]
  [Indiretta]
]
\end{forest}
\end{center}
\begin{itemize}
  \item \textbf{ossificazione diretta} (intramembranosa o mesenchimale): si forma direttamente da un tessuto connettivo primario, da cellule mesenchimali che si differenziano in osteoblasti;
  \item \textbf{ossificazione indiretta} (per sostituzione): si forma da una preesistente cartilagine, gradualmente sostituita da tessuto osseo.
\end{itemize}

% ---------------------------------------------------------------------------
\subsection{Ossificazione diretta}

\begin{enumerate}
  \item \textbf{sviluppo del centro di ossificazione}: centri di ossificazione costituiti da cellule stellate sparse in abbondante matrice extracellulare; le cellule mesenchimali si differenziano in osteoblasti, che secernono una matrice preossea (\textbf{osteoide}); gli osteoblasti si ordinano intorno all'osteoide in file di singole cellule (disposizione epitelioide);
  \item \textbf{osteociti che depositano sali minerali (mineralizzazione)}: successivamente l'osteoide acquisisce i sali minerali, diventando matrice mineralizzata o calcificata $\rightarrow$ prima trabecola ossea;
  \item \textbf{formazione delle trabecole}: la crescita della prima trabecola ossea avviene per apposizione di ulteriore tessuto osteoide da parte degli osteoblasti; gli osteoblasti restano inclusi nelle lacune scavate nella matrice ossea, emettono prolungamenti all'interno dei canalicoli e si differenziano in osteociti; altre trabecole, derivanti da centri di ossificazione vicini, confluiscono tra loro;
  \item \textbf{sviluppo del periostio, del tessuto osseo spugnoso e del tessuto osseo compatto}.
\end{enumerate}

% ---------------------------------------------------------------------------
\subsection{Ossificazione indiretta}

\rubrica{Ossificazione pericondrale}
L'osso si forma alla periferia dell'abbozzo cartilagineo. Dalla periferia si dipartono vasi in profondità; da questi, cellule si differenziano in condroclasti, che degradano la matrice cartilaginea.

\rubrica{Ossificazione endocondrale}
L'osso si forma all'interno dell'abbozzo cartilagineo, attraverso: proliferazione dei condrociti; ipertrofizzazione e apoptosi dei condrociti; allargamento e confluenza delle capsule cartilaginee; invasione di cellule mesenchimali; deposizione di tessuto osseo.

Al centro della diafisi i condrociti smettono di proliferare e ipertrofizzano; la matrice ialina calcifica per deposizione di sali di calcio. Le cellule pericondriali si differenziano in osteoblasti che depongono tessuto osseo (ossificazione diretta) $\rightarrow$ \textbf{manicotto osseo pericondrale}. Si formano vasi sanguigni dal pericondrio verso l'interno della matrice cartilaginea calcificata; le cellule cartilaginee ipertrofiche muoiono, lasciando lacune che si colmeranno di cellule osteoprogenitrici. La matrice amorfa cartilaginea viene degradata dai condroclasti; gli osteoblasti depongono tessuto osseo.

\rubrica{Fase 1}
Mentre la cartilagine si espande, i condrociti vicini al centro della diafisi aumentano notevolmente di dimensioni. La matrice viene ridotta ad una serie di piccole strutture, che presto cominceranno a calcificare. I condrociti ingranditi (ipertrofici) in seguito muoiono e si disintegrano, lasciando cavità all'interno della cartilagine.

\rubrica{Fase 2}
Si formano vasi sanguigni intorno ai margini della cartilagine e le cellule pericondriali si trasformano in osteoblasti. La diafisi cartilaginea viene poi rivestita da uno strato osseo superficiale (\textbf{collare osseo}).

\rubrica{Fase 3}
I vasi sanguigni penetrano nella cartilagine ed invadono la regione centrale. Le cellule mesenchimali che migrano attraverso i vasi sanguigni si differenziano in osteoblasti, che iniziano a produrre osso spugnoso nel \textbf{centro di ossificazione primario} (diafisi). La formazione dell'osso si estende poi lungo la diafisi verso entrambe le estremità.

\rubrica{Fase 4}
Man mano che l'osso si accresce, avviene il rimodellamento ad opera degli osteoclasti, che formano la cavità midollare. L'osso diafisario diviene più spesso e la cartilagine vicina a ciascuna epifisi è sostituita da osso. L'ulteriore crescita conduce ad un aumento della lunghezza (fase 5) e del diametro dell'osso.

\rubrica{Fase 5}
Capillari ed osteoblasti migrano nell'epifisi, formando un \textbf{centro di ossificazione secondario}.

% ---------------------------------------------------------------------------
\subsection{Riparazione di una frattura}

\begin{enumerate}
  \item frattura e sanguinamento all'interno dell'area, con formazione di un \textbf{ematoma da frattura};
  \item formazione di un \textbf{callo interno} e di una rete di osso spugnoso, che unisce le superfici interne; un \textbf{callo esterno} di cartilagine ed osso stabilizza i margini esterni;
  \item la cartilagine del callo esterno viene sostituita da osso e colonne di osso spugnoso uniscono le estremità della frattura; frammenti di osso necrotico e le aree di osso vicine alla frattura vengono rimossi e sostituiti;
  \item un rigonfiamento segna la zona della frattura, che successivamente verrà rimodellata e rimarrà solo un lieve segno della frattura.
\end{enumerate}
